\section{引言}\label{sec:intro}

部分子分布函数（PDFs）是刻画强子高能结构的关键物理量。它们定义为无限动量强子内夸克与胶子的动量分布，并为描述强子—强子或轻子—强子对撞机的实验数据提供必要输入。然而，从第一性原理计算 PDFs 一直是强子物理学中长期悬而未决的问题，因为 PDFs 是以光锥关联定义的强子低能性质。传统格点 QCD 方法聚焦于计算其矩，但由于不同矩算符之间的幂次发散混合，只能确定最低的几个矩~\cite{Martinelli:1987zd,Martinelli:1988xs,Detmold:2001dv,Dolgov:2002zm}。然而重构 PDFs 需要其全部矩的信息。

过去几年里，人们提出了一种绕开上述困难的新方法，如今已被表述为大动量有效理论（LaMET）~\cite{Ji:2013dva,Ji:2014gla}。LaMET 为普遍地研究强子的部分子性质提供了可能。对于给定的部分子可观测量（如 PDFs），可以把光锥关联解构，使所得的量——称为准可观测量——不再具有坐标系不变性，但却与时间无关。准可观测量在无穷洛伦兹助推下趋于原来的部分子可观测量，并且按构造具有与之相同的红外（IR）行为。在动量有限但很大的强子中，准可观测量可以因子化为部分子可观测量与硬系数的卷积，并带有被强子动量压低的幂次修正~\cite{Ji:2013dva,Ji:2014gla,Xiong:2013bka,Ma:2017pxb,Izubuchi:2018srq,Ji:2015qla,Xiong:2015nua}。
自 LaMET 提出以来，无论是对该形式体系的理论理解，还是从格点 QCD 直接计算 PDFs，都取得了长足进展（参见近期综述~\cite{Cichy:2018mum} 及其中文献）。格点计算主要集中于同位旋矢量夸克 PDFs，因为它们不与胶子 PDFs 混合。尽管体积有限且格距相对较粗，在物理点从格点数据确定的最新核子同位旋矢量夸克 PDFs 已与从实验数据抽取的现象学结果表现出合理的一致性~\cite{Chen:2018xof,Lin:2018qky,Alexandrou:2018pbm,Alexandrou:2018eet,Liu:2018hxv,Dulat:2015mca,Ball:2017nwa,Harland-Lang:2014zoa,Nocera:2014gqa,Ethier:2017zbq}。

胶子部分子物理的格点计算更为关键，因为胶子在大型强子对撞机（LHC）以及美国未来的电子—离子对撞机（EIC）的实验中扮演着极其重要的角色。然而，
与利用 LaMET 抽取夸克 PDFs 的密集工作相比，针对胶子准PDF的研究所投入的力量要少得多。
在文献~\cite{Wang:2017qyg,Wang:2017eel} 中，人们首次尝试理解胶子准PDF的重正化性质，发现胶子准PDF会受到理论对称性所允许的其他胶子算符的混合污染。另一方面，在首次探索性计算胶子非极化 PDF 的工作~\cite{Fan:2018dxu} 中，作者假定其所用格点计算中的胶子准PDF是乘法重正化的。

在本通讯中，我们对定义胶子准PDF的规范不变的非定域算符的重正化进行系统研究，重点考察其幂次发散结构。沿着涉及我们之中两位作者的一项早期工作~\cite{Ji:2017oey}（另见~\cite{Ishikawa:2017faj,Green:2017xeu}），我们引入一个一般框架，把伴随表示中的威尔逊线替换为两个辅助“重夸克”场的乘积。可以证明，这样一个带辅助“重夸克”的理论在微扰论所有阶都是可重正化的。
随后我们将展示如何构造可乘法重正化的胶子准PDF算符，并说明如何在格点上非微扰地计算重正化因子。该分析还推广到极化情形。最后，我们给出重正化胶子准PDF的因子化公式，由此可以提取胶子 PDF。我们的结果扫除了通过格点模拟定义连续胶子准PDF的障碍。
